\section{Mixed canonical resets and prescribed source resets at exits}
\label{sec:mixedexits}
The no-reset continuation of Section~\ref{sec:conservative} is now
supplemented by the active-reset proof of A096 and the control and
dummy substitutions of A097 and A073. Their exact source versions
are recorded in the accompanying claim ledger. The statement below
also permits a specified reset at the root of a retained exit branch.
It retains the same original ancestor quota.
The source locators are A096, run
\path{20261021T210221Z-fleet03-A096-live-gap-ledger},
\path{DERIVATION.md}, Sections 1--3; A097, run
\path{20261021T232544Z-fleet03-A097-cross-route-transfer},
\path{DERIVATION.md}, Sections 1--7; and A073, run
\path{20261022T211534Z-fleet03-A073-certified-logarithms},
\path{DERIVATION.md}, Sections 1--4. The exit-source conversion
below is an additional scoped argument, pending independent review.

Let $a\ge1$, $b\ge0$, $r\ge1$, $n=a+b$, $d=b+r\ge2$,
$s=2^{-a}$, $t=2^{-n}$, $e=2^{1-n-r}$, and $\mu=1-2s$.
On active, dummy-data and control bits respectively, set
\[
 g=1-2A_a,\qquad F=1-2A_a+2A_{a+b+r}.
\]
Here $A_k$ denotes the all-zero point event on its specified $k$ bits.
Supply a normalized reference $R_{\rm raw}$ with finite source cost
$c_{\rm raw}\le C_0 I(g)$. Calibrate its empty mass to $\mu^2$,
obtaining $R$ and $c=\CE(P_g,R)\le c_{\rm raw}$.
The canonical complete floor has data marginal
\[
 R_e(T)=(1-e/s)R(T)+(e/s)\mu\,1_{T=\varnothing}+2et
\]
and the actual conditional control spectrum. Its complete cost is
at most $W=c+10e(a+b+r)$, by the supplied-source floor lemma above.

\begin{lemma}[The three local substitutions]
At a proper active query ($a\ge2$), marginalize and calibrate the
source for the on-child, and retain the complementary inherited
reference. Then $B_j(F)+(k+W_{\rm on})/2\le W-e$.
At $a=1$, use the explicit point spectrum at the final on-child;
the full substitution costs at most $W$.
At a dummy query ($b\ge1$), or a control query ($r\ge2$), retain
the inherited off-reference and reset the on-child to its specified
canonical floor. Then
\[
 B_j(F)+(c_{\rm off}+W_{\rm on})/2<W+8e.
 \tag{MX1}
\]
\end{lemma}
\begin{proof}
For the active query the exact nonvacuum erasure difference from
$g$ is $-2e(1-4s+e)\le0$. The vacuum pair transfers amplitude
from $2s$ to $\mu\ge2s$, which decreases
$\phi(x^2,y^2)$ when $x\ge y$. To check the latter assertion,
put $z=x/y\ge1$: $L(z)=\ln(1+z^2)-z\ln(1+z^{-2})$ vanishes
at one and has
$L'(z)=2(z+1)/(1+z^2)-\ln(1+z^{-2})>0$.
Indeed $\ln(1+z^{-2})\le z^{-2}$ and the resulting numerator
is $2z^3+z^2-1>0$. Thus $B_j(F)\le B_j(g)$.

Write $v=R(\{j\})$ and
$\kappa=(2s-e)^2/(4s^2-4se+2et)$. The complementary constant
child has inherited vacuum mass
$q_*(v)=(\mu+e)^2+\kappa[(1-e/s)v+2et]$.
The ratio $(\mu^2+v)/q_*(v)$ increases with $v$ and is maximized
at $v=1-\mu^2$. At this endpoint, $1-q_*\le6e$ and, writing
$h=e/s$, the exact bound is
\[
 q_*\ge1-5h/4+13h^2/16
       =27/52+(13/16)(h-10/13)^2>1/2.
\]
For completeness, $1-\kappa\le e/s$ gives
$q_*\ge1-sh(6-4h)+s^2h(4-3h)$; its minimum over
$0<s\le1/4$ is at $s=1/4$ because its difference from that
endpoint is $(1/4-s)h[(6-4h)-(s+1/4)(4-3h)]\ge0$.
The bracket is at least $4-5h/2>0$. The signed complementary
excess over the source constant-child cost is therefore less
than $18e$, using $\ln2>2/3$. Source-reference disintegration and
empty calibration pay $B_j(g)$ and both source-child costs from
$c$. Since $e_{\rm on}=2e$ and the dimension falls by one,
the floor surcharge releases $10e$, against less than $9e$ of
complementary excess. This proves the proper active inequality.
At $a=1$ the on-child is $-1+2A_{b+r}$, whose explicit true
spectrum has vacuum $(1-2^{1-b-r})^2$ and each nonempty mass
$2^{2-2(b+r)}$. Resetting to this law costs minus the inherited
KL. The complementary reference is retained, so (NR1) proves
the boundary without the false $a=1$ erasure comparison.

For either non-active query, exact coefficient counting gives
$\sum|\alpha\beta|=e(4-4s+2e-2e/s)\le4e$.
The inequality $\ln(1+z^2)\le z$ gives
$B_j(F)\le8e/\ln2<12e$. On an occupied nonempty source
label, the off-child inherited likelihood is at least
$(1-h)R(T)/(1+h/2)$, $h=e/s\le1/2$.
For a control the multiplier is
$[(2s-e)^2+e^2]/[4s^2-4se+2te]$; for a dummy it is
$(2s-e)^2/[4s^2-4se+2te]$, with an additional nonnegative
marginal term retained. Their reciprocals are at most $1+h/2$:
use $e\le t$, $t\le s$ and
$(2s-e)^2+e^2\ge4s(s-e)>0$ for the control, and
$t\le s/2$ with $(2s-e)^2\ge2s^2$ for the dummy.
The vacuum is no worse than $R(0)=\mu^2$. Thus
\[
 c_{\rm off}-c\le4s[2h+h/(2\ln2)]<11e.
\]
For a dummy, $V=R_{-j}$ has $V(0)\ge\mu^2$ and $V(0)<1$;
the latter follows from the positive source singleton mass.
Empty calibration gives
$\bar R(T)=V(T)(1-\mu^2)/(1-V(0))\ge R(T)$ on occupied
nonempty labels and leaves the occupied vacuum unchanged.
Consequently $c_{\rm on}\le c$. For a control the source $R$
is unchanged. Combining the erasure bound, both complete
terminals and $W_{\rm on}\le c+20e(a+b+r-1)$ yields a surcharge
less than $(12+11/2-10)e<8e$, proving (MX1).
\end{proof}

\begin{proposition}[A prescribed exit-source reset is paid]
At any non-active exit created by these mixed histories, before
any further query has been executed inside that exit branch,
one may replace its inherited reference by the following specified
reference: at a control exit use the current source $R$, with point
mass on the unused controls; at a dummy exit use the current source
marginal calibrated to $\mu^2$, again with point mass on controls.
The root source cost of this new exit reference is at most $c$.
One may then execute any finite physical restriction tree using
only source-reference marginalization and empty calibration, with
the explicit point reference on constant exits. All executed
erasures and all terminal costs of this source tree sum to at most
that same $c$. A positive signed exit reset is allowed; its cost
is paid by the original envelope, not assumed nonpositive.
\end{proposition}
\begin{proof}
The control exit is exactly $g$ with unused dummy coordinates.
The proposed normalized reference has complete cost $c$ and
supports its true law. At a dummy exit, the pointwise calibration
argument above bounds its cost by $c$. References may have zero
mass on controls because the true exit spectrum is concentrated
there. No occupied atom is removed.

For the subsequent pure-source tree, (NR1) pays every query by
the incoming complete source cost. Every calibration decreases
cost by the binary KL of its old and prescribed vacuum masses;
the conditional nonempty law stays fixed. At a constant child,
the explicit point reference decreases cost to zero. Finite
telescoping therefore pays the entire pure-source continuation
from its initial cost. This uses an explicitly inherited source
reference, not an oracle for arbitrary Boolean spectra.

At the original non-active parent, replacing the full off-cost
in (MX1) by this complete source-tree ledger costs at most $c$.
The substitution is now bounded by $W+2e$, since its erasure
is less than $12e$ and the floor surcharge releases $10e$.
In particular it still fits the old $W+8e$ charge, even when the
signed reset from the old off-reference is positive. Main-branch
debits already executed before the reset are unchanged.
The exit may be retained without queries while other prescribed
main-branch steps are booked, then converted once: the same final
ledger comparison retains every intervening debit. This does not
authorize a fresh reset after arbitrary prior queries inside the
exit or a reset to an unrelated source reference.
\end{proof}

\begin{theorem}[One original quota, all eligible live mixed orders]
For any finite adaptive live-branch history of active, dummy and
eligible control steps, with either inherited off-references or
the specified exit-source resets and continuations above, the
complete erasure-plus-terminal ledger is at most
\[
 W_{\rm root}+8(d-1)e_{\rm root}
 \le (2C_0+10)I(F_{\rm root})-12e_{\rm root}.
 \tag{MX2}
\]
If every non-active exit uses its specified source reset, the
coefficient 8 can be replaced by 2 and the final reserve by 18.
An independent untouched terminal keeps its full additive cost.
\end{theorem}
\begin{proof}
At every nonterminal live on-step, physical reach halves and
$e$ doubles. Thus $\rho e=e_{\rm root}$. There are at most
$d-1$ non-active steps before $b=0,r=1$. Proper active queries
release capacity; the final active boundary never increases its
envelope. Weighted substitution proves the first bound in (MX2)
for every stopping cut. At a retained exit use its old full bill;
at a converted exit use its entire source-tree bill. No other
terminal is changed or omitted.

Put $I_g=2as$, $E=ae$ and $X=de$. The exact physical influence
is $I(F)=I_g+X-E$, with $E/I_g=2^{-d}\le1/4$. For
$K=2C_0+10$ the initial reserve obeys
\[
 KI(F)-W\ge(C_0+10)(I_g-2E)+2C_0X\ge5I_g
              =5a2^d e.
\]
For integers $d\ge2$, $5\cdot2^d-8(d-1)\ge12$; equality
holds at two and the increment is $5\cdot2^d-8>0$.
For the all-source-exit policy use instead
$5\cdot2^d-2(d-1)\ge18$, with the same argument.
Independent terminal cross entropy adds with coefficient one
because the exit/source updates act only on the edited factor.
Thus its entropy or KL funds none of these resets.
\end{proof}

This closes only the specified exit-root transaction and its
source continuation. Arbitrary incoming profiles, a different
prescribed reference after queries already executed in an exit,
and a cheap initial reference for every Boolean source remain
unpaid. Extra channel and code bills are outside (MX2).
The full conjecture and both integrated-chain audits remain open.
